%=====================================================================

\section{Bounded reduction}\label{sec:bddred}

The purpose of this section is to remove the boundedness hypothesis from
the sharp Saint--Venant stability inequality. The input is the
\emph{bounded} estimate proved in the earlier part of the manuscript
(Theorem~\ref{thm:boundedSV}), which we restate here in energy form.

\begin{theorem}[Bounded sharp Saint--Venant]\label{thm:SV-bounded}
For every $n\ge2$ and every $R\ge2$ there is a constant
$c_{\mathrm{SV}}(n,R)>0$ such that
\begin{equation}\label{eq:SV-bounded}
   E(\Om)-E(B_1)\;\ge\;c_{\mathrm{SV}}(n,R)\,\Fr(\Om)^2
\end{equation}
for every open set $\Om\subset B_R$ with $\abs{\Om}=\omega_n$.
\end{theorem}

This is Theorem~\ref{thm:boundedSV}, restated in energy form: with the
constant $C(n,R,\rstar)$ there one may take
$c_{\mathrm{SV}}(n,R)=1/C(n,R,\rstar)$ and
$\dT(\Om)=E(\Om)-E(B_1)$, and the parameter $\rstar$ is fixed at, say,
$\rstar=\tfrac34$.

Our goal is the global statement.

\begin{theorem}[Global sharp Saint--Venant stability]\label{thm:SV-global}
There is a dimensional constant $c_{\mathrm{SV}}^{\mathrm{glob}}(n)>0$
such that
\begin{equation}\label{eq:SV-global}
   E(\Om)-E(B_1)\;\ge\;c_{\mathrm{SV}}^{\mathrm{glob}}(n)\,\Fr(\Om)^2
\end{equation}
for every open set $\Om\subset\R^n$ with $\abs{\Om}=\omega_n$.
Equivalently, after rescaling, for every open set $\Om$ of finite
measure, with $B^*$ the ball of the same volume,
\begin{equation}\label{eq:T-global}
   T(B^*)-T(\Om)\;=\;2\bigl(E(\Om)-E(B^*)\bigr)\;\ge\;
   2\,c_{\mathrm{SV}}^{\mathrm{glob}}(n)\,
   \Bigl(\tfrac{\abs{\Om}}{\omega_n}\Bigr)^{(n+2)/n}\,\Fr(\Om)^2 .
\end{equation}
\end{theorem}

The strategy is the constructive reduction of Brasco, De~Philippis and
Velichkov \cite[\S5]{BDV}. From a set $\Om$ with small Saint--Venant
deficit we manufacture a \emph{bounded} surrogate $\Otilde$, of the same
volume, of dimensionally bounded diameter, with comparable deficit and
with controlled loss of asymmetry; applying
Theorem~\ref{thm:SV-bounded} to $\Otilde$ at a fixed dimensional radius
and transferring the estimate back to $\Om$ closes the small-deficit
regime. The complementary large-deficit regime is trivial because
$\Fr<2$. The whole construction is explicit: no compactness, no
minimising sequence and no contradiction argument enter.

The two technical ingredients are an $L^\infty$ tail estimate for the
torsion function (Lemma~\ref{lem:Linfty}, proved by a De~Giorgi
iteration) and the truncation construction itself
(Lemma~\ref{lem:truncation}). Both are proved here in full rather than
cited.

\subsection{Standing notation and the standard tools we cite}
\label{ssec:notation}

Throughout, $n\ge2$ and all sets are open subsets of $\R^n$ of finite
Lebesgue measure. We write $u_\Om\in H^1_0(\Om)$ for the torsion
function, i.e. the unique weak solution of $-\Delta u_\Om=1$ in $\Om$,
$u_\Om=0$ on $\partial\Om$, equivalently the minimiser of
\begin{equation}\label{eq:energy-def}
   J_\Om(v)\;=\;\tfrac12\int_\Om\abs{\nabla v}^2\,dx-\int_\Om v\,dx,
   \qquad v\in H^1_0(\Om),
\end{equation}
and $E(\Om)=J_\Om(u_\Om)=\min_{H^1_0(\Om)}J_\Om$. Testing the equation
with $u_\Om$ gives the standard identities
\begin{equation}\label{eq:energy-identities}
   \int_\Om\abs{\nabla u_\Om}^2\,dx=\int_\Om u_\Om\,dx=T(\Om),
   \qquad
   E(\Om)=-\tfrac12\int_\Om u_\Om\,dx=-\tfrac12\,T(\Om).
\end{equation}
We use $\omega_n=\abs{B_1}$, $T_n=T(B_1)=\omega_n/(n(n+2))$, the Fraenkel
asymmetry $\Fr(\Om)=\inf_{x_0}\abs{\Om\,\triangle\,(B_1+x_0)}/\omega_n$
(for $\abs{\Om}=\omega_n$ this lies in $[0,2)$), and the scale-invariant
deficit
\begin{equation}\label{eq:deficit-def}
   D(\Om)\;=\;E(\Om)\abs{\Om}^{-(n+2)/n}-E(B_1)\abs{B_1}^{-(n+2)/n},
\end{equation}
so that $D(\Om)=\omega_n^{-(n+2)/n}\bigl(E(\Om)-E(B_1)\bigr)\ge0$ when
$\abs{\Om}=\omega_n$. The non-negativity of $D$ is the
\emph{Saint--Venant inequality} (equivalently Talenti's comparison
theorem): among sets of given volume the ball minimises $E$, i.e.
\begin{equation}\label{eq:saint-venant}
   E(\Om)\,\abs{\Om}^{-(n+2)/n}\;\ge\;E(B)\,\abs{B}^{-(n+2)/n}
   \qquad\text{for every ball }B.
\end{equation}
We treat \eqref{eq:saint-venant} as a known classical result.

The external analytic facts we invoke without proof are:
\begin{itemize}
\item the \emph{Sobolev inequality}
   $\norm{w}_{L^{2^*}(\R^n)}\le S_n\norm{\nabla w}_{L^2(\R^n)}$ for
   $w\in H^1(\R^n)$ when $n\ge3$ (with $2^*=2n/(n-2)$), and its $n=2$
   surrogate explained in Remark~\ref{rem:n2};
\item the \emph{Saint--Venant/Talenti inequality}~\eqref{eq:saint-venant},
   used in the truncation argument and in the global $L^\infty$ bound
   \eqref{eq:global-Linfty} below;
\item Theorem~\ref{thm:SV-bounded}, the bounded stability input;
\item the \emph{suboptimal} torsional stability bound
   $\Fr(\Om)^4\le C(n)D(\Om)$ (Lemma~\ref{lem:seed}), used only to seed
   the truncation iteration. This is the single nontrivial result we cite
   rather than prove; see Remark~\ref{rem:seed-status}.
\end{itemize}
Everything else --- in particular the fast-convergence iteration lemma
(Lemma~\ref{lem:DG-iteration2}), the De~Giorgi $L^\infty$ tail estimate
(Lemma~\ref{lem:Linfty}), the energy-comparison inequality, the truncation
iteration (Lemma~\ref{lem:truncation}) and the gluing --- is proved from
scratch.

We will also need the elementary pointwise sup-bound on the torsion
function. By Talenti's symmetrisation comparison the torsion function of
$\Om$ is dominated, after Schwarz symmetrisation, by that of the ball of
the same volume; in particular
\begin{equation}\label{eq:global-Linfty}
   \norm{u_\Om}_{L^\infty(\Om)}\;\le\;\norm{u_{B^*}}_{L^\infty(B^*)}
   \;=\;\frac{(\abs{\Om}/\omega_n)^{2/n}}{2n}
   \;=:\;\Lambda_n\,\abs{\Om}^{2/n},
   \qquad \Lambda_n:=\frac{\omega_n^{-2/n}}{2n},
\end{equation}
where $B^*$ is the ball with $\abs{B^*}=\abs{\Om}$ and we used
$u_{B_r}(x)=(r^2-\abs{x}^2)/(2n)$. For $\abs{\Om}=\omega_n$ this reads
$\norm{u_\Om}_{L^\infty(\Om)}\le 1/(2n)$. We record \eqref{eq:global-Linfty}
as a consequence of the cited Talenti comparison.

\subsection{The fast-convergence iteration lemma}\label{ssec:DG-lemma}

The De~Giorgi scheme reduces an $L^\infty$ bound to the vanishing of a
nonlinear recursion. We isolate the underlying numerical lemma and prove
it, so that nothing is cited at this step. The mechanism is that a
super-quadratic recursion with a small enough starting value beats any
fixed geometric growth factor $b^k$.

\begin{lemma}[Fast geometric convergence]\label{lem:DG-iteration2}
Let $(a_k)_{k\ge0}$ be non-negative reals with
\begin{equation}\label{eq:DG-recursion2}
   a_{k+1}\;\le\;C\,b^{\,k}\,a_k^{\,1+\beta},\qquad k\ge0,
\end{equation}
where $C>0$, $b>1$, $\beta>0$. If
\begin{equation}\label{eq:DG-smallness2}
   a_0\;\le\;C^{-1/\beta}\,b^{-1/\beta^2},
\end{equation}
then
\begin{equation}\label{eq:DG-conclusion}
   a_k\;\le\;b^{-k/\beta}\,a_0\qquad\text{for all }k\ge0,
\end{equation}
so $a_k\to0$.
\end{lemma}

\begin{proof}
Put $r:=b^{-1/\beta}\in(0,1)$ and $\theta:=C\,a_0^{\beta}$. The smallness
hypothesis \eqref{eq:DG-smallness2} is precisely $\theta\le r$. We prove
the target $a_k\le r^{k}a_0$ by induction; this is
\eqref{eq:DG-conclusion}, and $r<1$ then gives $a_k\to0$.

\emph{Base case} $k=0$: equality. \emph{Inductive step.} Assume
$a_k\le r^{k}a_0$. Using \eqref{eq:DG-recursion2},
\[
   a_{k+1}\le C\,b^{k}\,a_k^{1+\beta}
          =C\,b^{k}\,a_k\,a_k^{\beta}
          \le C\,b^{k}\,a_k\,\bigl(r^{k}a_0\bigr)^{\beta}
          =\bigl(C\,a_0^{\beta}\bigr)\,(b\,r^{\beta})^{k}\,a_k
          =\theta\,(b\,r^{\beta})^{k}\,a_k .
\]
By the definition of $r$ we have $r^{\beta}=b^{-1}$, hence $b\,r^{\beta}=1$
and the geometric factor collapses:
\[
   a_{k+1}\le\theta\,a_k\le r\,a_k\le r\cdot r^{k}a_0=r^{k+1}a_0,
\]
using $\theta\le r$ and the inductive hypothesis. This closes the
induction.
\end{proof}

\begin{remark}
The identity $b\,r^{\beta}=1$ with $r=b^{-1/\beta}$ is the algebraic
mechanism behind the convergence: the super-quadratic gain
$a_k^{1+\beta}$ converts, once $a_0$ is small, the growing factor $b^k$
into a contraction.
\end{remark}

\subsection{The De~Giorgi \texorpdfstring{$L^\infty$}{L-infinity} tail
estimate}\label{ssec:Linfty}

\begin{lemma}[$L^\infty$ tail estimate]\label{lem:Linfty}
Let $n\ge3$. There is a dimensional constant $C_\flat=C_\flat(n)>0$ such
that for every open $\Om$ with $\abs{\Om}=\omega_n$ and every $R\ge1$,
\begin{equation}\label{eq:Linfty}
   \norm{u_\Om}_{L^\infty(\Om\setminus B_{R+1})}
   \;\le\;C_\flat\,\abs{\Om\setminus B_R}^{1/n}.
\end{equation}
For $n=2$ the same holds with the exponent $1/n$ replaced by
$\tfrac12-\tfrac1p$ for any fixed $p\in(2,\infty)$
(Remark~\textup{\ref{rem:n2}}); this weaker but still super-linear power is
all that the truncation uses.
\end{lemma}

\begin{proof}
If $\abs{\Om\setminus B_R}=0$ then $u_\Om=0$ a.e.\ on
$\Om\setminus B_R\supset\Om\setminus B_{R+1}$ and \eqref{eq:Linfty} is
trivial, so assume $\abs{\Om\setminus B_R}>0$. We carry out a De~Giorgi
iteration on the annular region $\Om\setminus B_R$.

\smallskip
\emph{Set-up: shrinking radii, cut-offs, and rising levels.}
For $k\in\mathbb N_0$ put
\[
   R_k:=R+1-2^{-k},
\]
so that $R_0=R$, the radii increase to $R_\infty=R+1$, and
$R_k\ge R\ge1$. Let $\varphi_k\in C^{0,1}(\R^n)$ be the radial cut-off
\[
   \varphi_k(x):=\phi_k(\abs{x}),\qquad
   \phi_k=0\ \text{on}\ [0,R_{k-1}],\quad
   \phi_k=1\ \text{on}\ [R_k,\infty),\quad
   \phi_k\ \text{affine on}\ [R_{k-1},R_k],
\]
for $k\ge1$, and $\varphi_0\equiv1$ on $\{\abs{x}\ge R_0\}$, extended by
$0$ inside (we only use $\varphi_k$ for $k\ge1$ in the recursion). Then
$0\le\varphi_k\le1$, $\varphi_k\equiv0$ on $B_{R_{k-1}}$, $\varphi_k\equiv1$
on $\R^n\setminus B_{R_k}$, and
\[
   \abs{\nabla\varphi_k}\le\frac1{R_k-R_{k-1}}
   =\frac1{2^{-(k-1)}-2^{-k}}=2^{\,k}\le 2^{\,k+1}.
\]
Choose rising levels
\[
   s_k:=M\,\abs{\Om\setminus B_R}^{1/n}\,(1-2^{-k}),\qquad k\ge0,
\]
where $M=M(n)\ge1$ is a free constant fixed at the end. Thus $s_0=0$,
$s_k\uparrow s_\infty:=M\abs{\Om\setminus B_R}^{1/n}$, and
\begin{equation}\label{eq:s-diff}
   s_{k+1}-s_k=M\abs{\Om\setminus B_R}^{1/n}\,2^{-(k+1)} .
\end{equation}

\smallskip
\emph{Caccioppoli inequality from the truncated test function.}
Fix $k\ge1$. The function $(u_\Om-s_k)_+\in H^1_0(\Om)$ (it vanishes
where $u_\Om\le s_k$, in particular near $\partial\Om$ since $u_\Om\ge0$
and $s_k\ge0$), and $\varphi_k$ is bounded Lipschitz, so
\[
   v_k:=\varphi_k^{\,2}\,(u_\Om-s_k)_+\ \in\ H^1_0(\Om)
\]
is an admissible test function. Insert it into the weak formulation
$\int_\Om\nabla u_\Om\cdot\nabla v\,dx=\int_\Om v\,dx$. Using
$\nabla v_k=\varphi_k^2\nabla(u_\Om-s_k)_+
+2\varphi_k(u_\Om-s_k)_+\nabla\varphi_k$ and
$\nabla u_\Om\cdot\nabla(u_\Om-s_k)_+=\abs{\nabla(u_\Om-s_k)_+}^2$ (the
gradient of $u_\Om$ equals that of $(u_\Om-s_k)_+$ on $\{u_\Om>s_k\}$ and
both factors vanish elsewhere), we get
\[
   \int\varphi_k^2\abs{\nabla(u_\Om-s_k)_+}^2\,dx
   +2\int\varphi_k(u_\Om-s_k)_+\,\nabla\varphi_k\cdot\nabla(u_\Om-s_k)_+\,dx
   =\int\varphi_k^2(u_\Om-s_k)_+\,dx .
\]
The middle term is bounded by Young's inequality,
$2\abs{\varphi_k(u_\Om-s_k)_+\nabla\varphi_k\cdot\nabla(u_\Om-s_k)_+}
\le\tfrac12\varphi_k^2\abs{\nabla(u_\Om-s_k)_+}^2
+2\abs{\nabla\varphi_k}^2(u_\Om-s_k)_+^2$,
and absorbing the first piece on the left yields the Caccioppoli
estimate
\begin{equation}\label{eq:caccioppoli}
   \int\abs{\nabla\bigl(\varphi_k(u_\Om-s_k)_+\bigr)}^2\,dx
   \;\le\;C_1\!\int\abs{\nabla\varphi_k}^2(u_\Om-s_k)_+^2\,dx
        \;+\;C_1\!\int\varphi_k^2(u_\Om-s_k)_+\,dx,
\end{equation}
with an absolute constant $C_1$ (one checks $C_1=10$ works after expanding
$\abs{\nabla(\varphi_k w)}^2\le2\varphi_k^2\abs{\nabla w}^2
+2\abs{\nabla\varphi_k}^2 w^2$ and combining with the absorbed bound; the
precise value is immaterial). Here we wrote $w=(u_\Om-s_k)_+$.

\smallskip
\emph{Sobolev inequality.} Set
$w_k:=\varphi_k(u_\Om-s_k)_+\in H^1_0(\R^n)$ and
\[
   A_k:=\bigl\{\varphi_k(u_\Om-s_k)_+>0\bigr\}
       =\bigl\{u_\Om>s_k\bigr\}\cap\{\varphi_k>0\}
       \subset\{u_\Om>s_k\}\cap(\R^n\setminus B_{R_{k-1}}).
\]
For $n\ge3$ the Sobolev inequality and Hölder's inequality on the support
$A_k$ give
\begin{equation}\label{eq:sobolev}
   \int w_k^2\,dx
   \le\Bigl(\int w_k^{2^*}dx\Bigr)^{2/2^*}\abs{A_k}^{1-2/2^*}
   \le S_n^2\,\abs{A_k}^{2/n}\!\int\abs{\nabla w_k}^2\,dx,
\end{equation}
using $1-2/2^*=2/n$. (The $n=2$ case is handled in
Remark~\ref{rem:n2}; the resulting estimate has the same shape with $2/n$
replaced by any fixed exponent in $(0,1)$, which only changes the
dimensional constants.)

Combining \eqref{eq:sobolev} with \eqref{eq:caccioppoli},
$\abs{\nabla\varphi_k}\le2^{k+1}$ (so $\abs{\nabla\varphi_k}^2\le4^{k+1}$)
and $0\le\varphi_k\le1$, we obtain
\begin{equation}\label{eq:after-sobolev}
   \int w_k^2\,dx
   \le S_n^2\,C_1\,\abs{A_k}^{2/n}
   \Bigl(4^{k+1}\!\!\int_{A_k}\!(u_\Om-s_k)_+^2\,dx
        +\!\!\int_{A_k}\!(u_\Om-s_k)_+\,dx\Bigr),
\end{equation}
where for the second term we used
$\int\varphi_k^2(u_\Om-s_k)_+\le\int_{A_k}(u_\Om-s_k)_+$. By the global
sup-bound \eqref{eq:global-Linfty} at $\abs{\Om}=\omega_n$ we have
$0\le(u_\Om-s_k)_+\le u_\Om\le1/(2n)\le1$ on $\Om$, hence both
$(u_\Om-s_k)_+^2\le1$ and $(u_\Om-s_k)_+\le1$ there, so each integral over
$A_k$ is at most $\abs{A_k}$. Therefore
\begin{equation}\label{eq:after-sobolev2}
   \int w_k^2\,dx
   \le S_n^2\,C_1\,\bigl(4^{k+1}+1\bigr)\,\abs{A_k}^{2/n}\,\abs{A_k}
   \le C_2\,4^{k}\,\abs{A_k}^{1+2/n},
\end{equation}
with $C_2=C_2(n):=8\,S_n^2 C_1$ (using $4^{k+1}+1\le8\cdot4^{k}$).

\smallskip
\emph{From energy to a measure recursion.} Restrict the left-hand side of
\eqref{eq:after-sobolev2} to the smaller set where the \emph{next} level
is exceeded. On $\{u_\Om>s_{k+1}\}\cap(\R^n\setminus B_{R_{k}})$ we have
$\varphi_k=1$ (since this set lies outside $B_{R_k}$) and
$u_\Om-s_k>s_{k+1}-s_k$, hence
$w_k^2=(u_\Om-s_k)_+^2\ge(s_{k+1}-s_k)^2$ there. Therefore
\begin{equation}\label{eq:chebyshev}
   (s_{k+1}-s_k)^2\,
   \bigl|\{u_\Om>s_{k+1}\}\cap(\R^n\setminus B_{R_{k}})\bigr|
   \;\le\;\int_{\R^n}w_k^2\,dx .
\end{equation}
Introduce the (relative) super-level measures
\[
   \mu_k:=\bigl|\{u_\Om>s_k\}\cap(\R^n\setminus B_{R_{k-1}})\bigr|,
   \qquad k\ge1,
\]
which decrease in $k$ (both $s_k$ and $R_{k-1}$ increase). Since
$A_k\subset\{u_\Om>s_k\}\cap(\R^n\setminus B_{R_{k-1}})$ we have
$\abs{A_k}\le\mu_k$, and the left-hand set in \eqref{eq:chebyshev} is
exactly $\{u_\Om>s_{k+1}\}\cap(\R^n\setminus B_{R_{(k+1)-1}})$, of
measure $\mu_{k+1}$. Combining \eqref{eq:chebyshev},
\eqref{eq:after-sobolev2} and $\abs{A_k}\le\mu_k$,
\[
   (s_{k+1}-s_k)^2\,\mu_{k+1}
   \;\le\;C_2\,4^{k}\,\mu_k^{1+2/n}.
\]
Insert the gap \eqref{eq:s-diff},
$(s_{k+1}-s_k)^2=M^2\abs{\Om\setminus B_R}^{2/n}4^{-(k+1)}$:
\begin{equation}\label{eq:mu-recursion}
   \mu_{k+1}\;\le\;
   \frac{C_2\,4^{k}}{M^2\abs{\Om\setminus B_R}^{2/n}4^{-(k+1)}}\,
   \mu_k^{1+2/n}
   \;=\;\frac{4\,C_2}{M^2\abs{\Om\setminus B_R}^{2/n}}\,16^{\,k}\,
   \mu_k^{1+2/n}.
\end{equation}

\smallskip
\emph{Normalisation and the iteration lemma.} Set
$a_k:=\mu_k/\abs{\Om\setminus B_R}$. Since
$\mu_k\le\abs{\R^n\setminus B_{R_{k-1}}\cap\{u_\Om>s_k\}}
\le\abs{\Om\setminus B_R}$ (because $R_{k-1}\ge R_0=R$ and
$\{u_\Om>s_k\}\subset\Om$), we have $0\le a_k\le1$; in particular
$a_1\le1$. Dividing \eqref{eq:mu-recursion} by
$\abs{\Om\setminus B_R}$ and using
$\mu_k^{1+2/n}=\abs{\Om\setminus B_R}^{1+2/n}a_k^{1+2/n}$,
\[
   a_{k+1}
   \le\frac{4C_2}{M^2\abs{\Om\setminus B_R}^{2/n}}\,16^k\,
   \frac{\abs{\Om\setminus B_R}^{1+2/n}}{\abs{\Om\setminus B_R}}
   \,a_k^{1+2/n}
   =\frac{4C_2}{M^2}\,16^{k}\,a_k^{1+2/n}.
\]
valid for all $k\ge1$. To put this in the exact form
\eqref{eq:DG-recursion2} (indexed from $0$) we shift: set
$\widetilde a_j:=a_{j+1}$ for $j\ge0$. Then for every $j\ge0$,
\[
   \widetilde a_{j+1}=a_{j+2}
   \le\frac{4C_2}{M^2}\,16^{\,j+1}\,a_{j+1}^{1+2/n}
   =\underbrace{\frac{64\,C_2}{M^2}}_{=:C}\,16^{\,j}\,
     \widetilde a_j^{\,1+2/n},
\]
which is \eqref{eq:DG-recursion2} with
\[
   C:=\frac{64\,C_2}{M^2},\qquad b:=16,\qquad\beta:=\frac2n .
\]
Choose $M=M(n)$ so large that the smallness condition
\eqref{eq:DG-smallness2} holds at the starting value
$\widetilde a_0=a_1\le1$. Since $a_1\le1$ it suffices to require
\[
   1\;\le\;C^{-1/\beta}\,b^{-1/\beta^2}
   =\Bigl(\tfrac{64C_2}{M^2}\Bigr)^{-n/2}16^{-n^2/4},
   \quad\text{i.e.}\quad
   M^2\;\ge\;64\,C_2\,16^{\,n^2/4\cdot(2/n)}=64\,C_2\,16^{\,n/2},
\]
which holds for $M:=8\,(C_2)^{1/2}4^{\,n/2}$ (a dimensional constant,
since $C_2=C_2(n)$). With this $M$, Lemma~\ref{lem:DG-iteration2} gives
$\widetilde a_j\to0$, i.e.\ $a_k\to0$ and hence $\mu_k\to0$.

\smallskip
\emph{Conclusion.} Letting $k\to\infty$ in $\mu_k\to0$ and using
$s_k\uparrow s_\infty=M\abs{\Om\setminus B_R}^{1/n}$ and
$R_{k-1}\uparrow R+1$, monotone convergence yields
\[
   \bigl|\{u_\Om>s_\infty\}\cap(\R^n\setminus B_{R+1})\bigr|
   =\lim_{k\to\infty}\mu_k=0 .
\]
Hence $u_\Om\le s_\infty=M\abs{\Om\setminus B_R}^{1/n}$ a.e.\ on
$\Om\setminus B_{R+1}$. Since $u_\Om\ge0$, this is
\eqref{eq:Linfty} with $C_\flat:=M=8\,(C_2)^{1/2}4^{n/2}$.
\end{proof}

\begin{remark}[The case $n=2$]\label{rem:n2}
For $n=2$ the embedding $H^1_0\hookrightarrow L^{2^*}$ degenerates, but
the Gagliardo--Nirenberg--Sobolev inequality gives, for any fixed
$p\in(2,\infty)$,
$\norm{w}_{L^{p}(\R^2)}\le S_{2,p}\norm{\nabla w}_{L^2(\R^2)}$ for
$w\in H^1_0(\R^2)$ with finite-measure support. Repeating the Hölder step with
exponent $p$ replaces $2/n=1$ by $1-2/p\in(0,1)$, so
\eqref{eq:mu-recursion} becomes
$\mu_{k+1}\le C\,16^k\,\mu_k^{1+\beta}$ with $\beta:=1-2/p>0$, and
Lemma~\ref{lem:DG-iteration2} applies. Tracking the level scale: with
$s_\infty=M\abs{\Om\setminus B_R}^{\alpha}$, self-consistency of the
recursion forces $\alpha=\tfrac{p-2}{2p}=\tfrac12-\tfrac1p$ (for $n\ge3$
the analogous bookkeeping gives $\alpha=1/n$, and the two agree only there).
Hence, for $n=2$, the $L^\infty$ tail estimate holds in the slightly weaker
form
\[
   \norm{u_\Om}_{L^\infty(\Om\setminus B_{R+1})}
   \;\le\;C(p)\,\abs{\Om\setminus B_R}^{\,\frac12-\frac1p},
   \qquad\text{any fixed }p\in(2,\infty),
\]
in place of the exponent $1/n$ in \eqref{eq:Linfty}. Fixing once and for
all $p=4$, so that $\tfrac12-\tfrac1p=\tfrac14>0$, this is a genuine
super-linear power of the tail mass, which is all the truncation of
\S\ref{ssec:truncation} uses: the energy-comparison
\eqref{eq:energy-comp} and the tail recursion \eqref{eq:bk2-bound} hold
with the exponent $1+\tfrac1n$ replaced by $1+\kappa$, $\kappa=\tfrac14$,
and the geometric iteration \eqref{eq:bk2-self}--\eqref{eq:bodd-decay} goes
through for any $\kappa>0$ with $n$-dependent constants. Consequently the
global Theorem~\ref{thm:SV-global} holds for all $n\ge2$; only the
auxiliary exponent in \eqref{eq:Linfty} (and the resulting dimensional
constants) changes at $n=2$.
\end{remark}

\subsection{A suboptimal stability seed}\label{ssec:seed}

The truncation iteration needs one quantitative input that the De~Giorgi
machinery does not by itself supply: a crude (far from sharp) control of the
asymmetry by the deficit, used only to make the tail mass small at the first
radius and to bound the number of truncation steps. We prove it here from the
level-set identity already established, so the present section introduces no
external stability input.

\begin{lemma}[Suboptimal stability seed]\label{lem:seed}
There is a dimensional constant $C_\dagger=C_\dagger(n)>0$ such that every
open $\Om\subset\R^n$ of finite measure satisfies
\begin{equation}\label{eq:seed}
   D(\Om)\;\ge\;\frac{1}{C_\dagger}\,\Fr(\Om)^{4}.
\end{equation}
Equivalently, $\Fr(\Om)\le\bigl(C_\dagger D(\Om)\bigr)^{1/4}$.
\end{lemma}

\begin{proof}
Since $D$ and $\Fr$ are scale invariant we may assume $\abs\Om=\omega_n$, so
that $D(\Om)=\omega_n^{-(n+2)/n}\dT(\Om)$ and $B^*=B_1$; write $A:=\Fr(\Om)\in
[0,2)$ and recall $E_\rho=\{u>t(\rho)\}$, $\abs{E_\rho}=\omega_n\rho^n$,
$a(\rho)=-t'(\rho)$. We use three facts proved earlier in the manuscript,
each of which holds for every open set of finite measure (their proofs do not
use the confinement $\Om\subset B_R$): the exact deficit identity
(Theorem~\ref{thm:id-main}); the profile-gap identity
\eqref{eq:profile-gap-rho}, namely $0\le a(\rho)\le\rho/n$ and
$\int_0^1\rho^n(\tfrac\rho n-a(\rho))\,d\rho=\tfrac2{\omega_n}\dT(\Om)$; and
the superlevel transfer (Lemma~\ref{lem:superlevel-transfer}),
$A\le\Fr(E_\rho)+2(1-\rho^n)$.

\emph{Step 1 (a level-set lower bound for $D_I$).}
Fix a finite-perimeter level $\rho$ and apply the constructive Fraenkel
quantitative isoperimetric inequality (Theorem~\ref{thm:constructive-fraenkel})
to $E_\rho$.  It gives a centre $z$ such that
$\abs{E_\rho\Delta B_\rho(z)}^2\le C_{\rm iso}(n)\,\rho^{\,n+1}
(\Per(E_\rho)-\Per(B_\rho))$. Since
$\Fr(E_\rho)\le\abs{E_\rho\Delta B_\rho(z)}/\abs{E_\rho}$ and
$\Per(E_\rho)-\Per(B_\rho)\le D_I(t(\rho))/\Per(B_\rho)$, with
$\abs{E_\rho}=\omega_n\rho^n$ and $\Per(B_\rho)=n\omega_n\rho^{n-1}$, this
rearranges to
\begin{equation}\label{eq:seed-iso}
   D_I(t(\rho))\;\ge\;c_1(n)\,\rho^{\,2n-2}\,\Fr(E_\rho)^2,
   \qquad c_1(n):=\frac{n\,\omega_n^{3}}{C_{\rm iso}(n)} .
\end{equation}

\emph{Step 2 (restrict to the outer collar).}
Put $\rho_A:=(1-A/4)^{1/n}$. For $\rho\in[\rho_A,1]$ one has
$1-\rho^n\le1-\rho_A^n=A/4$, so the superlevel transfer gives
$\Fr(E_\rho)\ge A-2(1-\rho^n)\ge A/2$, and $\rho^n\ge\rho_A^n=1-A/4>1/2$.
Drop $D_H\ge0$ in Theorem~\ref{thm:id-main} and change variables
$t=t(\rho)$ ($dt=a(\rho)\,d\rho$, $m=\omega_n\rho^n$); keeping only
$\rho\in[\rho_A,1]$, where the integrand is nonnegative, and inserting
\eqref{eq:seed-iso},
\[
   2\dT(\Om)\ \ge\
   \int_{\rho_A}^1\frac{D_I(t(\rho))}{n^2\omega_n^{2/n}(\omega_n\rho^n)^{1-2/n}}
       \,a(\rho)\,d\rho
   \ \ge\ c_2(n)\int_{\rho_A}^1\rho^{\,n}\,\Fr(E_\rho)^2\,a(\rho)\,d\rho,
\]
where $c_2(n)=c_1(n)/(n^2\omega_n)$ (using
$(\omega_n\rho^n)^{1-2/n}=\omega_n^{1-2/n}\rho^{n-2}$ and
$\rho^{2n-2}/\rho^{n-2}=\rho^n$). With $\Fr(E_\rho)\ge A/2$,
\begin{equation}\label{eq:seed-collar}
   2\dT(\Om)\ \ge\ \frac{c_2(n)\,A^2}{4}\int_{\rho_A}^1\rho^n a(\rho)\,d\rho .
\end{equation}

\emph{Step 3 (the profile gap bounds the collar mass below).}
By the profile-gap identity and $\tfrac\rho n-a\ge0$,
\[
   \int_{\rho_A}^1\rho^n a\,d\rho
   =\int_{\rho_A}^1\frac{\rho^{n+1}}{n}\,d\rho
     -\int_{\rho_A}^1\rho^n\Bigl(\frac\rho n-a\Bigr)d\rho
   \ \ge\ \frac{1-\rho_A^{\,n+2}}{n(n+2)}-\frac2{\omega_n}\dT(\Om).
\]
Since $\rho_A\le1$, $\rho_A^{\,n+2}\le\rho_A^n=1-A/4$, so
$1-\rho_A^{\,n+2}\ge A/4$ and
\begin{equation}\label{eq:seed-mass}
   \int_{\rho_A}^1\rho^n a\,d\rho\ \ge\ \frac{A}{4n(n+2)}-\frac2{\omega_n}\dT(\Om).
\end{equation}

\emph{Step 4 (conclusion).} Insert \eqref{eq:seed-mass} into
\eqref{eq:seed-collar}:
\[
   2\dT(\Om)\ \ge\ \frac{c_2(n)A^2}{4}
   \Bigl(\frac{A}{4n(n+2)}-\frac2{\omega_n}\dT(\Om)\Bigr)
   = c_3(n)\,A^3-c_4(n)\,A^2\,\dT(\Om),
\]
with $c_3(n)=c_2(n)/(16n(n+2))$ and $c_4(n)=c_2(n)/(2\omega_n)$. Hence
$\dT(\Om)\,(2+c_4(n)A^2)\ge c_3(n)A^3$, and since $A<2$,
\[
   \dT(\Om)\ \ge\ \frac{c_3(n)}{2+4c_4(n)}\,A^3\ =:\ c_5(n)\,A^3 .
\]
Finally $A<2$ gives $A^4\le2A^3$, so $\dT(\Om)\ge\tfrac12 c_5(n)A^4$;
multiplying by $\omega_n^{-(n+2)/n}$ gives \eqref{eq:seed} with
$C_\dagger(n)=2\,\omega_n^{(n+2)/n}/c_5(n)$.
\end{proof}

\begin{remark}[The seed is internal]\label{rem:seed-status}
The only non-internal ingredient in this proof is the constructive Fraenkel
quantitative isoperimetric inequality (Theorem~\ref{thm:constructive-fraenkel});
everything else -- the deficit identity, the profile-gap identity, and the
superlevel transfer -- is proved here. Thus the bounded reduction does not use
the strong-form input and depends on no sharp stability estimate external to the
paper. The exponent $4$ (indeed the $3$ obtained in the proof) is far
from sharp -- the sharp exponent $2$ is the content of the main theorem --
but only a fixed positive power of $D$ is used below. A suboptimal torsional
stability bound of this kind is classical; cf.\ Fusco--Maggi--Pratelli
\cite{FMP}.
\end{remark}

\subsection{The truncation construction}\label{ssec:truncation}

We now build the bounded surrogate. The estimate driving the construction
is an \emph{energy-comparison inequality}: cutting $\Om$ off at radius
$\sim k$ costs only a higher-order amount of energy, controlled by the
tail mass through Lemma~\ref{lem:Linfty}.

\begin{lemma}[Bounded truncation]\label{lem:truncation}
There exist dimensional constants $C_\sharp=C_\sharp(n)>0$,
$\delta=\delta(n)\in(0,1]$ and $d=d(n)\ge1$ such that for every open
$\Om$ with $\abs{\Om}=\omega_n$ and $D(\Om)\le\delta$ there is an open set
$\Otilde$ with
\begin{align}
   \abs{\Otilde}\;&=\;\omega_n,\label{eq:trunc-vol}\\
   \diam(\Otilde)\;&\le\;d,\label{eq:trunc-diam}\\
   D(\Otilde)\;&\le\;C_\sharp\,D(\Om),\label{eq:trunc-D}\\
   \Fr(\Om)\;&\le\;\Fr(\Otilde)+C_\sharp\,D(\Om).\label{eq:trunc-A}
\end{align}
\end{lemma}

\begin{proof}
Normalise so that the ball realising the asymmetry of $\Om$ is $B_1$
(translate $\Om$; both $\Fr$ and $D$ are translation invariant), i.e.
$\Fr(\Om)=\abs{\Om\,\triangle\,B_1}/\omega_n$. We will repeatedly use the
tail-mass notation
\begin{equation}\label{eq:bk-def}
   b_k:=\frac{\abs{\Om\setminus B_k}}{\omega_n},\qquad k\ge1,
\end{equation}
which is non-increasing in $k$ and satisfies $0\le b_k\le1$.

\smallskip
\emph{Step 0: the tail mass is small.}
Since $B_1$ realises the asymmetry,
$\abs{\Om\setminus B_1}\le\abs{\Om\,\triangle\,B_1}=\omega_n\Fr(\Om)$.
Combined with the suboptimal seed Lemma~\ref{lem:seed},
\begin{equation}\label{eq:b1-small}
   b_1\;\le\;\Fr(\Om)\;\le\;\bigl(C_\dagger\,D(\Om)\bigr)^{1/4}.
\end{equation}
In particular, since $b_k\le b_1$ for all $k\ge1$, every tail mass is
$\le(C_\dagger D(\Om))^{1/4}$, which can be made smaller than any
prescribed dimensional threshold by shrinking $\delta(n)$ (recall
$D(\Om)\le\delta(n)$). This is the only use of Lemma~\ref{lem:seed} in
the construction.

\smallskip
\emph{Step 1: the energy-comparison inequality.}
For $k\ge1$ let $\varphi_k$ be the Lipschitz cut-off
\[
   \varphi_k(x):=\min\bigl\{1,(k+2-\abs{x})_+\bigr\},
\]
so $\varphi_k\equiv1$ on $B_{k+1}$, $\varphi_k\equiv0$ outside $B_{k+2}$,
$0\le\varphi_k\le1$ and $\abs{\nabla\varphi_k}\le1$ everywhere (it is
$1$ only on the annulus $B_{k+2}\setminus B_{k+1}$). Then
$u_k:=\varphi_k\,u_\Om\in H^1_0(\Om\cap B_{k+2})$ is admissible for the
variational problem defining $E(\Om\cap B_{k+2})$, so
$E(\Om\cap B_{k+2})\le J_{\Om\cap B_{k+2}}(u_k)
=\tfrac12\int\abs{\nabla u_k}^2-\int u_k$. We use the pointwise algebraic
identity
\[
   \abs{\nabla(\varphi_k u_\Om)}^2
   =\nabla u_\Om\cdot\nabla(u_\Om\varphi_k^2)
   +\abs{\nabla\varphi_k}^2u_\Om^2,
\]
which is verified by expanding both sides
($\nabla(u_\Om\varphi_k^2)=\varphi_k^2\nabla u_\Om
+2\varphi_k u_\Om\nabla\varphi_k$, so the right-hand side equals
$\varphi_k^2\abs{\nabla u_\Om}^2+2\varphi_k u_\Om\nabla u_\Om\cdot
\nabla\varphi_k+\abs{\nabla\varphi_k}^2 u_\Om^2
=\abs{\varphi_k\nabla u_\Om+u_\Om\nabla\varphi_k}^2$). Integrating and
applying the weak equation
$\int\nabla u_\Om\cdot\nabla(u_\Om\varphi_k^2)=\int u_\Om\varphi_k^2$
with the admissible test function $u_\Om\varphi_k^2\in H^1_0(\Om)$, we
get
\[
   J_{\Om\cap B_{k+2}}(u_k)
   =\tfrac12\!\int u_\Om\varphi_k^2
   +\tfrac12\!\int\abs{\nabla\varphi_k}^2u_\Om^2
   -\int\varphi_k u_\Om .
\]
Now use $\tfrac12\varphi_k^2-\varphi_k=-\tfrac12(1-(1-\varphi_k)^2)
=-\tfrac12+\tfrac12(1-\varphi_k)^2$ to write
$\tfrac12\int u_\Om\varphi_k^2-\int\varphi_k u_\Om
=-\tfrac12\int u_\Om+\tfrac12\int(1-\varphi_k)^2u_\Om
=E(\Om)+\tfrac12\int(1-\varphi_k)^2u_\Om$, where we used
$E(\Om)=-\tfrac12\int u_\Om$ from \eqref{eq:energy-identities}.
Therefore
\begin{equation}\label{eq:energy-exact}
   E(\Om\cap B_{k+2})\le J_{\Om\cap B_{k+2}}(u_k)
   =E(\Om)+\tfrac12\int(1-\varphi_k)^2u_\Om
          +\tfrac12\int\abs{\nabla\varphi_k}^2u_\Om^2 .
\end{equation}
Both correction integrals are supported in $\R^n\setminus B_{k+1}$:
indeed $1-\varphi_k=0$ on $B_{k+1}$ and $\nabla\varphi_k=0$ on
$B_{k+1}\cup(\R^n\setminus B_{k+2})$. Hence, with $0\le1-\varphi_k\le1$
and $\abs{\nabla\varphi_k}\le1$,
\begin{equation}\label{eq:energy-comp-pre}
   E(\Om\cap B_{k+2})-E(\Om)
   \;\le\;\tfrac12\,\abs{\Om\setminus B_{k+1}}
   \Bigl(\sup_{\Om\setminus B_{k+1}}u_\Om
        +\sup_{\Om\setminus B_{k+1}}u_\Om^2\Bigr).
\end{equation}
By Lemma~\ref{lem:Linfty} applied with $R=k\ge1$ (so $R+1=k+1$ and the
tail region is $\Om\setminus B_{k+1}$),
$\sup_{\Om\setminus B_{k+1}}u_\Om\le C_\flat\abs{\Om\setminus B_k}^{1/n}
=C_\flat\,\omega_n^{1/n}b_k^{1/n}$, and likewise the square is bounded by
$C_\flat^2\omega_n^{2/n}b_k^{2/n}$. Also
$\abs{\Om\setminus B_{k+1}}=\omega_n b_{k+1}\le\omega_n b_k$. Plugging in,
and using $b_k\le1$ so that $b_k^{2/n}\le b_k^{1/n}$,
\begin{equation}\label{eq:energy-comp}
   E(\Om\cap B_{k+2})\;\le\;E(\Om)+C\,b_{k+1}\bigl(b_k^{1/n}+b_k^{2/n}\bigr)
   \;\le\;E(\Om)+C\,b_k^{\,1+1/n},
\end{equation}
for a dimensional constant $C=C(n)$ (we absorbed
$\tfrac12\omega_n(C_\flat\omega_n^{1/n}+C_\flat^2\omega_n^{2/n})$ and
$b_{k+1}\le b_k$ into $C$). This is the energy-comparison inequality.

\smallskip
\emph{Step 2: turning energy comparison into a tail recursion.}
By the Saint--Venant inequality \eqref{eq:saint-venant} applied to the
set $\Om\cap B_{k+2}$, whose volume is
$\abs{\Om\cap B_{k+2}}=\omega_n(1-b_{k+2})$,
\[
   E(B_1)\abs{B_1}^{-(n+2)/n}
   \le E(\Om\cap B_{k+2})\,\abs{\Om\cap B_{k+2}}^{-(n+2)/n},
\]
that is
\[
   E(B_1)\,\omega_n^{-(n+2)/n}\bigl(\omega_n(1-b_{k+2})\bigr)^{(n+2)/n}
   \le E(\Om\cap B_{k+2}),
\]
i.e.
\[
   E(B_1)\,(1-b_{k+2})^{(n+2)/n}\le E(\Om\cap B_{k+2}).
\]
Combine with \eqref{eq:energy-comp} and the definition
$E(\Om)=E(B_1)+\omega_n^{(n+2)/n}D(\Om)$ (which is
\eqref{eq:deficit-def} at $\abs{\Om}=\omega_n$):
\[
   E(B_1)(1-b_{k+2})^{(n+2)/n}
   \le E(\Om)+C\,b_k^{1+1/n}
   = E(B_1)+\omega_n^{(n+2)/n}D(\Om)+C\,b_k^{1+1/n}.
\]
Rearranging, and recalling $E(B_1)<0$, divide by $-E(B_1)>0$:
\[
   1-(1-b_{k+2})^{(n+2)/n}
   \le\frac{\omega_n^{(n+2)/n}D(\Om)+C\,b_k^{1+1/n}}{-E(B_1)}
   =:\widehat C_1\bigl(D(\Om)+b_k^{1+1/n}\bigr),
\]
with $\widehat C_1=\widehat C_1(n)>0$
(here $\widehat C_1=\max\{\omega_n^{(n+2)/n},C\}/(-E(B_1))$). To extract
$b_{k+2}$ we use the elementary lower bound
\begin{equation}\label{eq:power-lb}
   1-(1-t)^{p}\;\ge\;t\qquad\text{for all }t\in[0,1],\ p\ge1,
\end{equation}
which holds because $r\mapsto(1-t)^{r}$ is non-increasing on $[0,\infty)$,
so $(1-t)^{p}\le(1-t)^{1}=1-t$ for $p=(n+2)/n\ge1$. Applying
\eqref{eq:power-lb} with $t=b_{k+2}\in[0,1]$ and $p=(n+2)/n$,
\begin{equation}\label{eq:bk2-bound}
   b_{k+2}\;\le\;1-(1-b_{k+2})^{(n+2)/n}
   \;\le\;\widehat C\,\bigl(D(\Om)+b_k^{\,1+1/n}\bigr),
   \qquad \widehat C:=\widehat C_1 .
\end{equation}

\smallskip
\emph{Step 3: the geometric iteration to a bounded stopping index.}
Define the stopping index
\begin{equation}\label{eq:K-def}
   K:=\max\bigl\{k\ge1:\ b_k\ge2\widehat C\,D(\Om)\bigr\},
\end{equation}
with the convention $K:=0$ if $b_1<2\widehat C D(\Om)$ (then the tail is
already controlled at the first radius and Step~5 applies with $K=0$).
The set in \eqref{eq:K-def} is finite because $b_k\to0$ as $k\to\infty$
(the tail mass of a finite-measure set vanishes), so $K$ is well defined.
We prove
\begin{equation}\label{eq:K-bound}
   K\;\le\;K(n)
\end{equation}
for a dimensional $K(n)$, provided $\delta(n)$ is small enough. Assume
$K\ge3$ (otherwise \eqref{eq:K-bound} is trivial).

\emph{Self-improving contraction.} For every $k$ with $k+2\le K$ we have,
by definition of $K$, that $b_{k+2}\ge2\widehat C D(\Om)$, i.e.\
$\widehat C D(\Om)\le\tfrac12 b_{k+2}$. Feeding this into
\eqref{eq:bk2-bound},
\[
   b_{k+2}\le\widehat C D(\Om)+\widehat C b_k^{1+1/n}
   \le\tfrac12 b_{k+2}+\widehat C b_k^{1+1/n},
\]
and absorbing $\tfrac12 b_{k+2}$ on the left,
\begin{equation}\label{eq:bk2-self}
   b_{k+2}\;\le\;2\widehat C\,b_k^{\,1+1/n}
   \qquad\text{whenever }k+2\le K .
\end{equation}

\emph{Smallness from the seed.} By Step~0, $b_1\le(C_\dagger D(\Om))^{1/4}
\le(C_\dagger\delta(n))^{1/4}$. Choose $\delta(n)$ small enough that
\begin{equation}\label{eq:seed-smallness}
   b_1\;\le\;(C_\dagger\,\delta(n))^{1/4}\;\le\;(2\widehat C)^{-2n}
   \;=:\;\eta_n\in(0,1).
\end{equation}
Then $2\widehat C\,b_k^{1/(2n)}\le2\widehat C\,\eta_n^{1/(2n)}=1$ for all
$k\ge1$ (as $b_k\le b_1\le\eta_n$), so \eqref{eq:bk2-self} self-improves
to
\begin{equation}\label{eq:bk2-power}
   b_{k+2}\;=\;\bigl(2\widehat C\,b_k^{1/(2n)}\bigr)\,b_k^{\,1+1/(2n)}
   \;\le\;b_k^{\,1+1/(2n)}\qquad\text{whenever }k+2\le K .
\end{equation}

\emph{Iterating along odd indices.} Apply \eqref{eq:bk2-power} repeatedly,
starting from $k=1$, as long as the index stays $\le K$. Writing the odd
indices $1,3,5,\dots$ and setting $\gamma:=1+\tfrac1{2n}>1$, induction
gives, for every $j\ge0$ with $2j+1\le K$,
\begin{equation}\label{eq:bodd-decay}
   b_{2j+1}\;\le\;b_1^{\,\gamma^{\,j}} .
\end{equation}
Indeed \eqref{eq:bodd-decay} holds at $j=0$, and if it holds at $j$ with
$2(j+1)+1=2j+3\le K$, then \eqref{eq:bk2-power} (at $k=2j+1$, legitimate
since $2j+3\le K$) gives
$b_{2j+3}\le b_{2j+1}^{\gamma}\le\bigl(b_1^{\gamma^{j}}\bigr)^{\gamma}
=b_1^{\gamma^{j+1}}$.

\emph{Bounding $K$.} Let $j_K:=\lfloor(K-1)/2\rfloor$, the largest $j$
with $2j+1\le K$. By the maximality of $K$ and \eqref{eq:bodd-decay},
\begin{equation}\label{eq:K-pinch}
   2\widehat C\,D(\Om)\;\le\;b_{2j_K+1}\;\le\;b_1^{\,\gamma^{\,j_K}}
   \;\le\;\bigl(C_\dagger D(\Om)\bigr)^{\gamma^{\,j_K}/4},
\end{equation}
using $b_1\le(C_\dagger D(\Om))^{1/4}$. Take logarithms (recall
$D(\Om)\le\delta(n)<1$, so $\log D(\Om)<0$); writing $L:=-\log D(\Om)>0$,
\eqref{eq:K-pinch} reads
\[
   \log(2\widehat C)-L
   \;\le\;\frac{\gamma^{\,j_K}}{4}\bigl(\log C_\dagger-L\bigr).
\]
For $\delta(n)$ small, $L=-\log D(\Om)\ge-\log\delta(n)$ is large; in
particular we may take $\delta(n)$ small enough that
$L\ge2\max\{\abs{\log(2\widehat C)},\abs{\log C_\dagger}\}$, whence
$\log(2\widehat C)-L\ge-\tfrac32 L$ and $\log C_\dagger-L\le-\tfrac12 L$.
The displayed inequality then forces
$-\tfrac32 L\le-\tfrac{\gamma^{\,j_K}}{8}L$, i.e.\
$\gamma^{\,j_K}\le12$. Since $\gamma=1+\tfrac1{2n}>1$ this bounds
$j_K\le\log12/\log\gamma=:j_*(n)$, hence
$K\le2j_K+2\le2j_*(n)+2=:K(n)$, a dimensional constant. This proves
\eqref{eq:K-bound}.

\emph{Step 4: tail at the stopping radius.} By the maximality
\eqref{eq:K-def}, $b_{K+1}<2\widehat C D(\Om)$; we record this for the next
step.

\smallskip
\emph{Step 5: the truncated set and its volume.}
By the definition \eqref{eq:K-def} of $K$ and $K\le K(n)$,
$b_{K+1}<2\widehat C D(\Om)$, hence
\begin{equation}\label{eq:tail-controlled}
   \abs{\Om\setminus B_{K+1}}=\omega_n b_{K+1}<2\widehat C\,\omega_n\,D(\Om).
\end{equation}
In fact we cut one radius further out to have room for the cut-off: by
monotonicity $b_{K+3}\le b_{K+1}<2\widehat C D(\Om)$, so
\begin{equation}\label{eq:vol-cut}
   \abs{\Om\cap B_{K+3}}=\omega_n(1-b_{K+3})\ge\omega_n(1-2\widehat C D(\Om)),
   \qquad
   \abs{\Om\setminus B_{K+3}}<2\widehat C\,\omega_n D(\Om).
\end{equation}
Define
\[
   r:=\Bigl(\frac{\abs{\Om\cap B_{K+3}}}{\omega_n}\Bigr)^{1/n}
   =(1-b_{K+3})^{1/n}\in\bigl[(1-2\widehat C D(\Om))^{1/n},\,1\bigr],
   \qquad
   \Otilde:=r^{-1}\,(\Om\cap B_{K+3}).
\]
Then $\abs{\Otilde}=r^{-n}\abs{\Om\cap B_{K+3}}=\omega_n$, which is
\eqref{eq:trunc-vol}. Since $\Om\cap B_{K+3}\subset B_{K+3}$ and
$r\ge(1-2\widehat C\delta(n))^{1/n}\ge\tfrac12$ for $\delta(n)$ small,
\[
   \diam(\Otilde)=r^{-1}\diam(\Om\cap B_{K+3})\le r^{-1}\cdot2(K+3)
   \le4\bigl(K(n)+3\bigr)=:d(n),
\]
giving \eqref{eq:trunc-diam}. We record for later the linearised bound
\begin{equation}\label{eq:r-near-1}
   0\le1-r^n=b_{K+3}<2\widehat C D(\Om),
   \qquad
   0\le1-r\le1-r^n<2\widehat C D(\Om),
\end{equation}
the middle inequality because $r\le1$ implies $1-r\le1-r^n$.

\smallskip
\emph{Step 6: deficit of the surrogate.}
Apply the energy comparison \eqref{eq:energy-comp} at $k=K+1$ (so the
cut radius is $k+2=K+3$):
\[
   E(\Om\cap B_{K+3})\le E(\Om)+C\,b_{K+1}^{\,1+1/n}
   \le E(\Om)+C\,(2\widehat C D(\Om))^{1+1/n}
   \le E(\Om)+C'\,D(\Om),
\]
where in the last step we used $D(\Om)\le\delta(n)\le1$ so that
$D(\Om)^{1+1/n}\le D(\Om)$, and absorbed constants into
$C'=C'(n)$. Since $E$ scales as
$E(t\Omega')=t^{n+2}E(\Omega')$ and $\Otilde=r^{-1}(\Om\cap B_{K+3})$,
\[
   E(\Otilde)=r^{-(n+2)}E(\Om\cap B_{K+3})
   \le r^{-(n+2)}\bigl(E(\Om)+C'D(\Om)\bigr).
\]
Because $\abs{\Otilde}=\omega_n$,
$D(\Otilde)=\omega_n^{-(n+2)/n}(E(\Otilde)-E(B_1))$. Using
$E(\Om)=E(B_1)+\omega_n^{(n+2)/n}D(\Om)$ and $r\le1$ (so
$r^{-(n+2)}\ge1$), with $E(B_1)<0$,
\begin{align*}
   E(\Otilde)-E(B_1)
   &\le r^{-(n+2)}\bigl(E(B_1)+\omega_n^{(n+2)/n}D(\Om)+C'D(\Om)\bigr)-E(B_1)\\
   &=\bigl(r^{-(n+2)}-1\bigr)E(B_1)
     +r^{-(n+2)}\bigl(\omega_n^{(n+2)/n}+C'\bigr)D(\Om).
\end{align*}
For the first term, $r^{-(n+2)}-1\ge0$ and $E(B_1)<0$, so
$(r^{-(n+2)}-1)E(B_1)\le0$ and we may bound it by
$(r^{-(n+2)}-1)|E(B_1)|$. Quantitatively, by \eqref{eq:r-near-1},
$r^{n}\ge1-2\widehat C D(\Om)$, so
$r^{-(n+2)}=r^{-n}\cdot r^{-2}\le(1-2\widehat C D(\Om))^{-(n+2)/n}
\le1+C''D(\Om)$ for $\delta(n)$ small (Bernoulli/Taylor bound, since
$2\widehat C D(\Om)\le2\widehat C\delta(n)\le\tfrac12$). Hence
$r^{-(n+2)}-1\le C''D(\Om)$ and $r^{-(n+2)}\le1+C''D(\Om)\le2$, giving
\[
   E(\Otilde)-E(B_1)
   \le C''D(\Om)\cdot\abs{E(B_1)}+2\bigl(\omega_n^{(n+2)/n}+C'\bigr)D(\Om)
   \le C_2'\,D(\Om),
\]
with $C_2'=C_2'(n)$. Dividing by $\omega_n^{(n+2)/n}$,
\begin{equation}\label{eq:D-tilde}
   D(\Otilde)=\omega_n^{-(n+2)/n}\bigl(E(\Otilde)-E(B_1)\bigr)
   \le\omega_n^{-(n+2)/n}C_2'\,D(\Om)=:C_3\,D(\Om),
\end{equation}
which is \eqref{eq:trunc-D} with $C_3=C_3(n)$.

\smallskip
\emph{Step 7: asymmetry loss.}
Let $B_1(x_0)$ realise the asymmetry of $\Otilde$, so
$\Fr(\Otilde)=\abs{\Otilde\,\triangle\,B_1(x_0)}/\omega_n$. We compare the
asymmetry of $\Om$ to that of $\Otilde$ through the chain of inclusions
relating $\Om$, $\Om\cap B_{K+3}$ and $\Otilde=r^{-1}(\Om\cap B_{K+3})$.
For any test ball $B_1(y)$,
\[
   \abs{\Om\,\triangle\,B_1(y)}
   \le\abs{\Om\,\triangle\,(\Om\cap B_{K+3})}
     +\abs{(\Om\cap B_{K+3})\,\triangle\,B_1(y)}.
\]
The first term is $\abs{\Om\setminus B_{K+3}}<2\widehat C\omega_n D(\Om)$
by \eqref{eq:vol-cut}. For the second, scale by $r$: writing
$E_r:=\Om\cap B_{K+3}=r\Otilde$ and $B_1(y)=r\,B_{1/r}(y/r)$,
\[
   \abs{(r\Otilde)\,\triangle\,(r\,B_{1/r}(y/r))}
   =r^n\,\abs{\Otilde\,\triangle\,B_{1/r}(y/r)}
   \le r^n\Bigl(\abs{\Otilde\,\triangle\,B_1(y/r)}
              +\abs{B_1(y/r)\,\triangle\,B_{1/r}(y/r)}\Bigr).
\]
The last term is the volume difference of concentric balls of radii $1$
and $1/r\ge1$:
$\abs{B_1\,\triangle\,B_{1/r}}=\omega_n(r^{-n}-1)$. Hence
\[
   r^n\abs{B_1(y/r)\,\triangle\,B_{1/r}(y/r)}
   =r^n\,\omega_n(r^{-n}-1)=\omega_n(1-r^n)
   <2\widehat C\,\omega_n D(\Om),
\]
using \eqref{eq:r-near-1}. Choosing $y:=r x_0$ so that $y/r=x_0$ is the
optimal centre for $\Otilde$, the middle term becomes
$r^n\abs{\Otilde\,\triangle\,B_1(x_0)}=r^n\omega_n\Fr(\Otilde)
\le\omega_n\Fr(\Otilde)$ (as $r\le1$). Collecting,
\[
   \abs{\Om\,\triangle\,B_1(rx_0)}
   \le2\widehat C\omega_n D(\Om)+\omega_n\Fr(\Otilde)+2\widehat C\omega_n D(\Om),
\]
and dividing by $\omega_n$ and taking the infimum over translates (which
only decreases the left side),
\begin{equation}\label{eq:trunc-A-pre}
   \Fr(\Om)\le\frac{\abs{\Om\,\triangle\,B_1(rx_0)}}{\omega_n}
   \le\Fr(\Otilde)+4\widehat C\,D(\Om).
\end{equation}
This is \eqref{eq:trunc-A} with the constant $4\widehat C$. Taking
$C_\sharp:=\max\{C_3,4\widehat C\}=C_\sharp(n)$ proves
\eqref{eq:trunc-D}--\eqref{eq:trunc-A} simultaneously, and we already have
\eqref{eq:trunc-vol}--\eqref{eq:trunc-diam}. This completes the proof.
\end{proof}

\begin{remark}[Translation into a fixed ball]\label{rem:translate}
The diameter bound \eqref{eq:trunc-diam} is translation invariant, so
after a translation we may assume $\Otilde\subset B_{d(n)}$. Both
$\Fr$ and $D$ are translation invariant, so
\eqref{eq:trunc-vol}--\eqref{eq:trunc-A} are preserved.
\end{remark}

\subsection{The far branch}\label{ssec:far}

\begin{lemma}[Large-deficit quadratic bound]\label{lem:far}
For every open $\Om$ with $\abs{\Om}=\omega_n$ and
$D(\Om)\ge\delta(n)/4$,
\begin{equation}\label{eq:far}
   E(\Om)-E(B_1)\;\ge\;\omega_n^{(n+2)/n}\,\frac{\delta(n)}{16}\,\Fr(\Om)^2 .
\end{equation}
\end{lemma}

\begin{proof}
Since $\abs{\Om}=\omega_n$ gives $\Fr(\Om)<2$, we have
$\Fr(\Om)^2/4<1$. Hence
\[
   D(\Om)\ge\frac{\delta(n)}{4}
   =\frac{\delta(n)}{4}\cdot1
   \ge\frac{\delta(n)}{4}\cdot\frac{\Fr(\Om)^2}{4}
   =\frac{\delta(n)}{16}\Fr(\Om)^2 .
\]
Multiplying by $\omega_n^{(n+2)/n}$ and using
$E(\Om)-E(B_1)=\omega_n^{(n+2)/n}D(\Om)$ gives \eqref{eq:far}.
\end{proof}

\subsection{The near branch}\label{ssec:near}

Fix the dimensional working radius
\begin{equation}\label{eq:Rstar}
   R_*(n):=\max\{d(n),2\},
\end{equation}
so $R_*\ge d(n)$ (the surrogates of Lemma~\ref{lem:truncation} lie in
$B_{R_*}$ after translation, by Remark~\ref{rem:translate}) and
$R_*\ge2$ (the hypothesis of Theorem~\ref{thm:SV-bounded}).

\begin{lemma}[Small-deficit transfer]\label{lem:near}
Set
\begin{equation}\label{eq:csv-tilde}
   \widetilde c_{\mathrm{SV}}(n)
   :=\omega_n^{-(n+2)/n}\,c_{\mathrm{SV}}\bigl(n,R_*(n)\bigr)>0 .
\end{equation}
For every open $\Om$ with $\abs{\Om}=\omega_n$ and
$D(\Om)\le\delta(n)$,
\begin{equation}\label{eq:near-Fr}
   \Fr(\Om)^2\;\le\;
   \Bigl(\frac{4\,C_\sharp(n)}{\widetilde c_{\mathrm{SV}}(n)}
        +4\,C_\sharp(n)^2\Bigr)\,D(\Om).
\end{equation}
\end{lemma}

\begin{proof}
Apply Lemma~\ref{lem:truncation} to get $\Otilde$ satisfying
\eqref{eq:trunc-vol}--\eqref{eq:trunc-A}; by Remark~\ref{rem:translate}
assume $\Otilde\subset B_{d(n)}\subset B_{R_*(n)}$. Theorem
\ref{thm:SV-bounded} at $R=R_*(n)$ gives
$E(\Otilde)-E(B_1)\ge c_{\mathrm{SV}}(n,R_*)\Fr(\Otilde)^2$; dividing by
$\omega_n^{(n+2)/n}$ and using $\abs{\Otilde}=\omega_n$,
\begin{equation}\label{eq:D-tilde-A}
   D(\Otilde)\;\ge\;\widetilde c_{\mathrm{SV}}(n)\,\Fr(\Otilde)^2 .
\end{equation}
Combining \eqref{eq:D-tilde-A} with \eqref{eq:trunc-D},
\[
   \Fr(\Otilde)^2\le\frac{D(\Otilde)}{\widetilde c_{\mathrm{SV}}}
   \le\frac{C_\sharp}{\widetilde c_{\mathrm{SV}}}\,D(\Om).
\]
Now use \eqref{eq:trunc-A}, the elementary
$(a+b)^2\le2a^2+2b^2$, and $D(\Om)\le\delta(n)\le1$ (so
$D(\Om)^2\le D(\Om)$):
\begin{align*}
   \Fr(\Om)^2
   &\le\bigl(\Fr(\Otilde)+C_\sharp D(\Om)\bigr)^2
    \le2\Fr(\Otilde)^2+2C_\sharp^2D(\Om)^2\\
   &\le\frac{2C_\sharp}{\widetilde c_{\mathrm{SV}}}D(\Om)+2C_\sharp^2D(\Om)
    \le\Bigl(\frac{4C_\sharp}{\widetilde c_{\mathrm{SV}}}+4C_\sharp^2\Bigr)D(\Om),
\end{align*}
which is \eqref{eq:near-Fr}.
\end{proof}

\begin{corollary}[Small-deficit quadratic bound]\label{cor:near}
With
\begin{equation}\label{eq:cnear}
   c_{\mathrm{near}}(n):=
   \frac{\omega_n^{(n+2)/n}}
        {\,4C_\sharp(n)\bigl(1/\widetilde c_{\mathrm{SV}}(n)+C_\sharp(n)\bigr)}>0,
\end{equation}
every open $\Om$ with $\abs{\Om}=\omega_n$ and $D(\Om)\le\delta(n)$
satisfies $E(\Om)-E(B_1)\ge c_{\mathrm{near}}(n)\,\Fr(\Om)^2$.
\end{corollary}

\begin{proof}
Rearrange \eqref{eq:near-Fr} to
$D(\Om)\ge\Fr(\Om)^2/\bigl(4C_\sharp/\widetilde c_{\mathrm{SV}}
+4C_\sharp^2\bigr)$ and multiply by $\omega_n^{(n+2)/n}$, using
$E(\Om)-E(B_1)=\omega_n^{(n+2)/n}D(\Om)$.
\end{proof}

\subsection{Proof of the global theorem}\label{ssec:global-proof}

\begin{proof}[Proof of Theorem~\ref{thm:SV-global}]
Set
\begin{equation}\label{eq:csv-glob}
   c_{\mathrm{SV}}^{\mathrm{glob}}(n)
   :=\min\Bigl\{\,c_{\mathrm{near}}(n),\;
   \omega_n^{(n+2)/n}\,\tfrac{\delta(n)}{16}\,\Bigr\}>0 .
\end{equation}
Let $\Om$ be open with $\abs{\Om}=\omega_n$.

\emph{Case A: $D(\Om)\ge\delta(n)/4$.} Lemma~\ref{lem:far} gives
$E(\Om)-E(B_1)\ge\omega_n^{(n+2)/n}(\delta(n)/16)\Fr(\Om)^2
\ge c_{\mathrm{SV}}^{\mathrm{glob}}(n)\Fr(\Om)^2$.

\emph{Case B: $D(\Om)\le\delta(n)$} (which includes
$D(\Om)\le\delta(n)/4$). Corollary~\ref{cor:near} gives
$E(\Om)-E(B_1)\ge c_{\mathrm{near}}(n)\Fr(\Om)^2
\ge c_{\mathrm{SV}}^{\mathrm{glob}}(n)\Fr(\Om)^2$.

The two cases overlap on $\delta(n)/4\le D(\Om)\le\delta(n)$ and together
cover $[0,\infty)$; in either case \eqref{eq:SV-global} holds with
$c_{\mathrm{SV}}^{\mathrm{glob}}(n)$ as in \eqref{eq:csv-glob}, a
dimensional constant because $\delta(n)$, $c_{\mathrm{near}}(n)$ depend
only on $n$ (recall $R_*(n)$ and hence $c_{\mathrm{SV}}(n,R_*(n))$ are
dimensional).

\emph{General volume.} For arbitrary finite-measure $\Om$ put
$\Om_*:=(\omega_n/\abs{\Om})^{1/n}\Om$, so $\abs{\Om_*}=\omega_n$,
$\Fr(\Om_*)=\Fr(\Om)$ and $D(\Om_*)=D(\Om)$ by scale invariance of
$\Fr$ and $D$. Applying \eqref{eq:SV-global} to $\Om_*$ reads
$D(\Om)\ge\omega_n^{-(n+2)/n}c_{\mathrm{SV}}^{\mathrm{glob}}(n)\Fr(\Om)^2$.
Multiplying by $\abs{\Om}^{(n+2)/n}$ and using
$E(\Om)\abs{\Om}^{-(n+2)/n}-E(B^*)\abs{B^*}^{-(n+2)/n}=D(\Om)$ with
$\abs{B^*}=\abs{\Om}$, then $E=-T/2$, gives
\[
   T(B^*)-T(\Om)=2\bigl(E(\Om)-E(B^*)\bigr)
   \ge2\,c_{\mathrm{SV}}^{\mathrm{glob}}(n)
   \Bigl(\tfrac{\abs{\Om}}{\omega_n}\Bigr)^{(n+2)/n}\Fr(\Om)^2,
\]
which is \eqref{eq:T-global}.
\end{proof}

\begin{remark}[Sharpness of the exponent]
The exponent $2$ on $\Fr$ in \eqref{eq:SV-global} is optimal: it is
saturated by smooth, volume-preserving ellipsoidal perturbations of
$B_1$, for which both $E(\Om)-E(B_1)$ and $\Fr(\Om)^2$ are of the same
quadratic order in the perturbation parameter.
\end{remark}


